\section{Independent checks across scales}

The common coefficients connect several additional comparisons. The
calculations below retain their own selections and covariance treatments;
overlapping catalogues and spectra are not multiplied into a single
significance. The execution register identifies completed diagnostics,
open likelihoods, and their source files. It is an exploratory local
record, rather than an external blind preregistration. \source{U1}

\subsection{Local galaxies and weak lensing}
For 130 selected SPARC galaxies with 2710 radial points, profiling the
stated stellar mass-to-light, distance, and inclination priors gives a
common-$a_0$ penalized score of 11258.0. Adding a
$5\,\kms$ velocity floor gives 3433.2, with median per-galaxy
$\chi^2$ per point of 1.54. The sensitivity to the adopted velocity
errors is substantial. This is a diagonal-error algebraic-RAR check;
it does not replace the full director-field calculation or account for
all correlated noncircular motions. For 100 flat-velocity galaxies,
the baryonic Tully--Fisher log residual has median 0.020 dex and rms
0.067 dex. \source{U1, S03}

The isolated-lens audit distinguishes the older KiDS release with its
covariance from the updated published table. At the reference $a_0$,
the latter gives a diagonal statistic of 10.06 for fifteen bins, or
5.99 for seven inner bins, including the tabulated statistical and
interpolation/extrapolation uncertainties. Its full table covariance
was unavailable, and its published $H_0=73$ convention was retained.
The two reductions overlap and provide no independent combined evidence.
\source{U1, S04}

\subsection{CMB spectra, lensing, and velocity comparisons}
The positive-allocation cosmology of Section~\ref{sec:cosmology} was
also evaluated against 135 ACT primary-spectrum bins, profiling the
stated calibration at fixed cosmological parameters. The score is
162.73 against 162.02 for the reference \LCDM\ fit. The ten-bin ACT
lensing comparison gives 16.21 against 17.86. A higher-accuracy
lensing calculation changes the bandpowers by 0.013\%. These are
separate checks; adding their scores to Planck would require the
appropriate overlap cuts and joint covariance. \source{U1, S05--S06}

The lensing--velocity statistic $E_G$ depends on the Weyl potential and
the baryon velocity response. The calculation gives 0.4375 at
$z=0.316$ and 0.3862 at $z=0.555$, nearly constant over
$k=0.005$--$0.15\,\mathrm{Mpc}^{-1}$. The cited ACT--BOSS estimates
are $0.49^{+0.14}_{-0.11}$ and $0.31^{+0.06}_{-0.05}$.
The no-slip identity alone does not determine this observable.
In the conserved-fluid implementation, CAMB's ordinary
\texttt{delta\_tot} omits part of the bookkeeping capacity component;
a physical total-matter $\sigma_8$ must be reconstructed independently
of the arbitrary reference-dust split. \source{U1, S09}

The best-fit vacuum density reaches half its present value by
$z=7.928$ and has $\rho_\Lambda(2)/\rho_\Lambda(0)=0.96824$.
Its effective equation of state today is $-1.00054$. A conserved-CDM
$w_0,w_a$ inference and this interacting release model must be compared
through the same distances and perturbations. The phantom sign alone
does not settle the DESI comparison. \source{C5}

\subsection{Stellar systems, constants, and tensor modes}
With the parent-carrier rule, a wide binary has no independent capacity
halo. A radial Galactic-tide estimate for one to two solar masses at
10--30 kAU gives $\Delta v/v\simeq1.1\times10^{-5}$--$6\times10^{-4}$.
The range depends on separation, mass, and orientation; a universal
$10^{-4}$ ceiling is unwarranted. The selected Solar quadrupole
$Q_2=0$ is 0.89 quoted standard deviations from the Cassini central
value. No full Gaia selection-and-contamination likelihood has been
performed. \source{U1, S10}

The selected particle-coefficient route gives
$G=6.6742890813\times10^{-11}\,\mathrm{m^3\,kg^{-1}\,s^{-2}}$,
$-0.073$ quoted CODATA standard deviations; the muon ratio is
$206.768280237$, at $-0.535$ quoted standard deviations. Its tau mass
is 1776.9188 MeV against $1776.93\pm0.09$ MeV. Of the 401 tested
routing integers, only the integer $7$ lies within two quoted standard deviations
for the stated route; of 43 tested Farey candidates, only $2/7$ lies
within one. These finite-menu uniqueness checks do not quantify the
historical search over possible formulas. The measured Newton scale
and marked-formula residuals were known during development.
\source{U1, S11}

For a constant Einstein coefficient and minimally coupled scalar matter,
the tensor quadratic action has two positive-energy luminal modes.
Public GWTC-4 polarization comparisons give combined
$\log_{10}\mathcal B_{S/T}=-14.72$ and
$\log_{10}\mathcal B_{V/T}=-5.33$ for pure scalar and vector
alternatives. Mixed tensor--scalar and tensor--vector alternatives
have values $-0.20$ and $+0.10$, respectively. This supports
compatibility with the minimal tensor branch and leaves mixed sectors
weakly constrained; it supplies no distinctive preference over GR.
The native tensor continuum and a map from higher-curvature operators
to propagation or Lorentz-violation bounds remain open.
\source{U1, S12; D1}

Three eight-Gyr halo phase realizations retain late stock tilts of
$39.58^\circ$--$42.36^\circ$ and stellar tilts of
$21.38^\circ$--$22.33^\circ$. Rounding and precession persist.
These are repeatable outcomes of the specified kinetic model, with
force convergence and stationary formation still unresolved. The
aligned-branch failure remains part of the comparison. The earlier
Crater II, Antlia 2, and tidal-dwarf calculations retain their stated
selection and dynamical limits. \source{U1, S13--S14; G3--G7}

